% 事件列的上极限与下极限（教学示意）。三个样本结果对前16个事件的归属：
% omega_1：恰在偶数项事件中（无穷多次，但并非最终总在）→ 属于上极限，不属于下极限；
% omega_2：只在第1,2,4,7项事件中 → 两者都不属于；
% omega_3：在第2,5项及第8项起的每个事件中 → 属于下极限，因而也属于上极限。
% 阴影格表示 omega 属于 A_n；右侧给出按所述规律延续时的结论。
\begin{tikzpicture}[line width=.5pt,font=\fontsize{8}{10}\selectfont,>=Stealth,
  cell/.style={draw=FigureGrid,minimum width=4.6mm,minimum height=5mm,inner sep=0pt},
  on/.style={cell,fill=FigureBlueFill,draw=FigureBlue},
  lab/.style={font=\fontsize{7.5}{9}\selectfont,text=FigureMuted}]
 \foreach \n in {1,...,16}{\node[text=FigureInk] at (4.6mm*\n,6.5mm) {$\n$};}
 \node[text=FigureInk,anchor=east] at (2mm,6.5mm) {$n$};
 % row 1: even n
 \foreach \n in {1,...,16}{
   \pgfmathparse{int(mod(\n,2))}
   \ifnum\pgfmathresult=0 \node[on] at (4.6mm*\n,0) {}; \else \node[cell] at (4.6mm*\n,0) {}; \fi}
 % row 2: n in {1,2,4,7}
 \foreach \n in {1,...,16}{\node[cell] at (4.6mm*\n,-6mm) {};}
 \foreach \n in {1,2,4,7}{\node[on] at (4.6mm*\n,-6mm) {};}
 % row 3: n in {2,5} and n>=8
 \foreach \n in {1,...,16}{\node[cell] at (4.6mm*\n,-12mm) {};}
 \foreach \n in {2,5,8,9,10,11,12,13,14,15,16}{\node[on] at (4.6mm*\n,-12mm) {};}
 \foreach \y/\w in {0/1,-6mm/2,-12mm/3}{
   \node[text=FigureInk,anchor=east] at (2mm,\y) {$\omega_{\w}$};
   \node[text=FigureMuted] at (4.6mm*17+1mm,\y) {$\cdots$};}
 \node[lab,anchor=west] at (4.6mm*18,0) {偶数项都出现};
 \node[lab,anchor=west] at (4.6mm*18,-6mm) {第7项以后不再出现};
 \node[lab,anchor=west] at (4.6mm*18,-12mm) {第8项起每项都出现};
 \node[text=FigureInk] at (4.6mm*18+33mm,6.5mm) {上极限};
 \node[text=FigureInk] at (4.6mm*18+44mm,6.5mm) {下极限};
 \foreach \y/\a/\b in {0/属于/不属于,-6mm/不属于/不属于,-12mm/属于/属于}{
   \node[text=FigureInk] at (4.6mm*18+33mm,\y) {\a};
   \node[text=FigureInk] at (4.6mm*18+44mm,\y) {\b};}
\end{tikzpicture}
